% Not All Errors Are Equal: Consequence-Aware Reasoning Compute Allocation
% Target: ICLR 2027 (deadline ~Sept-Oct 2026)

\documentclass{article}
\usepackage{iclr2026_conference,times}

\input{math_commands.tex}

\usepackage{hyperref}
\usepackage{url}

\title{Not All Errors Are Equal:\\Consequence-Aware Reasoning Compute Allocation}

% The \author macro works with any number of authors. Use \And or \AND
% to break the author list across columns; we go anonymous for submission.
\author{Anonymous authors\\
Paper under double-blind review}

\newcommand{\fix}{\marginpar{FIX}}
\newcommand{\new}{\marginpar{NEW}}

\iclrfinalcopy % Uncomment for camera-ready

\begin{document}

\maketitle

\begin{abstract}
Modern reasoning models---OpenAI's o-series, DeepSeek-R1, Claude with
extended thinking, and Qwen3-Thinking---allocate test-time compute
adaptively across tasks: tasks that look harder receive more thinking
tokens, while easier ones stop earlier. Recent work on adaptive
test-time compute codifies this difficulty-based heuristic into
learned routers and confidence-driven early-stopping policies. All of
these methods share an implicit assumption: compute buys accuracy,
and accuracy is fungible across tasks.

This assumption breaks in real deployment. The ``remaining error''
hidden behind a given accuracy carries wildly different costs across
tasks: a typo in a log message and a database migration that
silently corrupts a production table both count as a single unit of
accuracy on a benchmark, but only one of them is worth additional
compute to avoid. We argue that test-time compute should be
allocated by expected cost-of-error, not by predicted difficulty.

We propose \emph{consequence-aware} test-time compute allocation. At
its core is a lightweight predictor that estimates the
deployment-time cost-of-error of a task directly from its issue
text; a scheduler then routes higher-consequence tasks to larger
compute tiers---or larger thinking budgets---under the same total
budget as a difficulty-based baseline.

We evaluate on 700 software-engineering tasks (SWE-bench Lite and
Multi-SWE-bench mini), labeled under three independent pipelines, and
audit three contemporary thinking models. Across all labeling
pipelines, consequence is statistically orthogonal to difficulty and
remains a significant predictor of cost-weighted gain after
controlling for difficulty and cross-model confidence. Current
thinking models do not allocate compute by consequence: spending is
either uncorrelated with consequence, saturated at the token cap, or
weakly leaning but far below the required magnitude. Our issue-only
predictor never misclassifies a high-consequence task as
low-consequence in 300 SWE-bench tasks. Under matched compute
budgets, consequence-aware routing reduces cost-weighted loss by
more than $30\%$ relative to difficulty-aware routing, and the
predictor-driven variant retains over $90\%$ of the oracle gain.
Notably, difficulty-aware routing performs worse than random because
the hardest tasks are unsolvable at any tier.
\end{abstract}

\section{Introduction}

% TODO: vignette + formal objective + contributions

\section{Related Work}

% TODO: Snell 2024 / Damani 2024 / Manvi 2024 / Russell-Wefald / Yang-Lin

\section{Consequence as a Construct}

% TODO: definition + F1 orthogonality + P0b confounding

\section{Model Audit}

% TODO: Qwen3-8B / Qwen3-VL / Claude n=171

\section{Predicting Consequence at Deployment Time}

% TODO: predictor + class-2 to class-0 = 0 safety

\section{Method: Cost-Weighted Compute Allocation}

% TODO: formal objective + scheduling algorithm

\section{Experiments}

% TODO: Phase 3b + marginal utility + (#4 in appendix)

\section{Discussion and Limitations}

% TODO

\section{Conclusion}

% TODO

\bibliography{references}
\bibliographystyle{iclr2026_conference}

\end{document}
